% ===========================================================================
% FIGURE — the process / isolation model, and the problem it exists to solve.
% Verified against kayfabe-core/src/gpu.rs (Proc, Vas, Channel, IsolateBox),
% gpa.rs (GpaSpace/GpaArena), kayfabe-isolate/src/lib.rs (IsolateFactory).
% ===========================================================================
\begin{tikzpicture}[x=1mm,y=1mm, font=\sffamily\scriptsize]
\def\tc#1{{\ttfamily\fontsize{5.4}{6}\selectfont #1}}

% ---------------------------------------------------------------- the guest --
\node[guestb, minimum width=54mm, minimum height=14mm] (pa) at (-48,89)
  {\textbf{guest process A} \;(e.g.\ PyTorch)\\[2pt]
   \tiny \tc{hClient = 0xc1d00001}\\[-0.5pt]
   \tiny buffer at guest VA \tc{0x7f\_0000\_0000}\\[-0.5pt]
   \tiny channel handle \tc{0x00000010}};
\node[guestb, minimum width=54mm, minimum height=14mm] (pb) at (48,89)
  {\textbf{guest process B} \;(e.g.\ llama.cpp)\\[2pt]
   \tiny \tc{hClient = 0xc1d00001} \;\textbf{— the same value}\\[-0.5pt]
   \tiny buffer at guest VA \tc{0x7f\_0000\_0000} \;\textbf{— the same VA}\\[-0.5pt]
   \tiny channel handle \tc{0x00000010} \;\textbf{— the same handle}};
\node[align=center, font=\sffamily\tiny, color=kfhostln, text width=32mm] at (0,89)
  {\textbf{Not a corner case.}\\ Handles and guest VAs are\\ \emph{per-process namespaces};
   two\\ processes colliding on every\\ value is normal.\\[1.5pt]
   \textbf{This is bug \#14.}};

\draw[ar, kfguestln] (pa.south) -- ++(0,-4);
\draw[ar, kfguestln] (pb.south) -- ++(0,-4);

% --------------------------------------------------- what disambiguates them --
\node[emulb, minimum width=134mm, minimum height=10mm] at (0,74)
  {\textbf{The only values that are NOT process-chosen are the ones the guest \emph{kernel} assigns from real hardware:}\\[2pt]
   \tiny \tc{Pdb} — the page-directory base, which the GMMU itself keys page tables by
     $\Rightarrow$ \textbf{the memory boundary}\\[-0.5pt]
   \tiny \tc{VChid} — the virtual channel id a doorbell token decodes to
     $\Rightarrow$ \textbf{the execution boundary}\\[-0.5pt]
   \tiny Everything is keyed \tc{(GpuId, Pdb)} or \tc{(GpuId, VChid)}. Never on a client handle.
     Never on host CPU state — no CR3 is read anywhere.};

% ------------------------------------------------------------- the two Procs --
\foreach \x/\lbl/\pdb/\ch/\ar in {-46/A/0x0a1000/14/{[0x0000\_0000,\;0x1000\_0000)},
                                   46/B/0x0b7000/15/{[0x1000\_0000,\;0x2000\_0000)}}
{
  \draw[draw=kfcoreln, line width=1pt, rounded corners=2pt, fill=kfcore!30]
    (\x-31,11) rectangle (\x+31,65);
  \node[anchor=north, font=\sffamily\bfseries\small, color=kfcoreln] at (\x,64)
    {\tc{Proc} \lbl\; — one guest process};
  \node[coreb, minimum width=58mm, minimum height=10mm] at (\x,53)
    {\textbf{address plane}\\[1pt]
     \tiny \tc{Vas} per \tc{(GpuId, Pdb=\pdb)} — its own \tc{AddressTable}\\[-0.5pt]
     \tiny \textbf{and its own host address space} (the \#14 fix)};
  \node[coreb, minimum width=58mm, minimum height=10mm] at (\x,41)
    {\textbf{execution plane}\\[1pt]
     \tiny \tc{Channel} per \tc{vChid = \ch} + a per-proc \tc{ExecPlane}\\[-0.5pt]
     \tiny nothing scheduled is one-shot or device-global (the \#12 fix)};
  \node[coreb, minimum width=58mm, minimum height=10mm] at (\x,29)
    {\textbf{completion plane}\\[1pt]
     \tiny its own \tc{CompletionQueue} + \tc{FenceArms}; re-delivery is\\[-0.5pt]
     \tiny driven off \emph{this} proc's own poll (the starvation fix)};
  \node[coreb, minimum width=58mm, minimum height=10mm] at (\x,17)
    {\textbf{isolate + GPA arena}\\[1pt]
     \tiny a disjoint slice of the guest-physical window: \tc{\ar}\\[-0.5pt]
     \tiny \tc{GpaSpace::release} takes the arena \textbf{by value}};
}
\draw[arb, kfhostln, dashed, line width=0.8pt] (-16,41) -- (16,41);
\node[align=center, font=\sffamily\tiny, color=kfhostln, text width=30mm] at (0,35)
  {\textbf{no shared state}\\ on any of the\\ four planes};

% ------------------------------------------------------------- the isolates --
\node[portb, fill=white, minimum width=58mm, minimum height=11mm] (ia) at (-48,4)
  {\textbf{isolate process} for \tc{(Proc A, GPU 0)}\\[1pt]
   \tiny separate OS process, unprivileged uid, namespaces,\\[-0.5pt]
   \tiny \tc{pivot\_root}, seccomp, cleared env and fds\\[-0.5pt]
   \tiny a bounded pool of single-in-flight workers};
\node[portb, fill=white, minimum width=58mm, minimum height=11mm] (ib) at (48,4)
  {\textbf{isolate process} for \tc{(Proc B, GPU 0)}\\[1pt]
   \tiny its own host RM client; handles minted here are\\[-0.5pt]
   \tiny typed to \emph{this} isolate, so a cross-isolate handle\\[-0.5pt]
   \tiny does not merely fail — it does not type-check};
\draw[ar, kfcoreln] (-48,11) -- (ia.north);
\draw[ar, kfcoreln] ( 48,11) -- (ib.north);

\node[align=center, font=\sffamily\tiny, color=kfhostln, text width=32mm] at (0,4)
  {\textbf{Blast radius.} A guest process that\\ crashes, wedges or attacks its own\\ isolate loses that isolate
   and nothing\\ else. The claim is \textbf{process-grade},\\ not VM-grade, isolation between\\ guest processes.};

% -------------------------------------------------------------- the host GPU --
\node[hostb, minimum width=134mm, minimum height=10mm] (hg) at (0,-13)
  {\textbf{one real host GPU}, driven by the ordinary NVIDIA driver through ordinary unprivileged \tc{ioctl}s\\[2pt]
   \tiny the host's own kernel RM keeps the two isolates' address spaces, channels and contexts apart —
   we are not asking it for anything unusual};
\draw[ar] (ia.south) -- (ia.south |- hg.north);
\draw[ar] (ib.south) -- (ib.south |- hg.north);

% ---------------------------------------------------------------- honesty box --
\node[box, fill=kfhost!30, draw=kfhostln, minimum width=178mm, minimum height=11mm] at (0,-29)
  {\textbf{\color{kfhostln}What is actually running, as of this document (2026-08-20).}\\[2pt]
   \tiny The per-\tc{Proc} plane ownership above is \textbf{built, heavily tested against mocks, and
   now driven on real hardware} — \tc{kayfabe-isolate-host} is a real sandboxed, capability-less,\\[-0.5pt]
   \tiny six-namespace child holding \file{/dev/nvidiactl} and \file{/dev/nvidia0}, and it is what
   executes the guest's compute. Its \tc{RmBackend} issues real host RM ioctls.\\[-0.5pt]
   \tiny ★★★ \textbf{The multi-\tc{Proc} axis this figure exists to argue has been driven}: two
   CONCURRENT guest CUDA processes both computed (w299, 2026-08-14) — 3 isolates materialized,
   3 live, 0 refusing, distinct \tc{proc=} ids —\\[-0.5pt]
   \tiny and the C artifact's \#14 hang did NOT reproduce, even with the second started while the
   first was inside \tc{cuCtxCreate}. \textbf{\color{kfhostln}⊘ Two identical short workloads,
   emulated CE work, one VM.}\\[-0.5pt]
   \tiny \textbf{\color{kfhostln}A second context SEQUENTIALLY in one boot still hangs}, and
   \textbf{\color{kfhostln}there is no tenant axis at all} — no \tc{VmId} anywhere in
   \file{crates/}; cross-VM separation is incidental to two host processes.};

\end{tikzpicture}
